\section*{Exercises}
Exercises~1--4 are core; 5--6 are transfer. Full solutions appear in the solution part.
\begin{exercise}\label{ex:1.1}Evaluate $2(x+3)$ and $2x+3$ at $x=4$. Explain the difference in words.\end{exercise}
\begin{exercise}\label{ex:1.2}Prove that the sum of an even integer and an odd integer is odd. State the integer that witnesses the conclusion.\end{exercise}
\begin{exercise}\label{ex:1.3}An assistant claims that the product of two odd integers is odd because $(2k+1)^2$ is odd. Identify the coverage problem and give a complete proof.\end{exercise}
\begin{exercise}\label{ex:1.4}Let $a_1=2$, $a_2=-1$, and $a_3=4$. Expand and evaluate $\sum_{j=1}^3(2a_j+1)$.\end{exercise}
\begin{exercise}\label{ex:1.5}Find every real solution of $x(x-2)=0$. Explain why division by $x$ alone is not a complete method. You may use the arithmetic fact that a product of real numbers is zero only when at least one factor is zero.\end{exercise}
\begin{exercise}\label{ex:1.6}You check a numerical claim for inputs $1$ through $10$. Construct a formula that agrees with the claim $f(n)=0$ on all these inputs but fails at $11$. Hint: use a product with one factor for each checked input.\end{exercise}
